\documentclass[12pt,a4paper,reqno]{amsart}

% language
\usepackage[greek,english]{babel}
\usepackage[utf8]{inputenc}
\usepackage{alphabeta}

% change default names to greek
\addto\captionsenglish{
    \renewcommand{\contentsname}{Περιεχόμενα}
    \renewcommand{\refname}{Βιβλιογραφία}
    \renewcommand{\datename}{Ημερομηνία:}
    \renewcommand{\urladdrname}{Ιστοσελίδα}
}

% math 
\usepackage{amsmath,amsthm,amssymb,amscd}

% font
\usepackage[cal=euler]{mathalfa}
\usepackage{libertinus-type1}
% \usepackage{txfonts} % for upright greek letters
\usepackage{bm} % for bold symbols
\usepackage{bbm} % for the simply-looking bb symbols

% miscellaneous 
\usepackage{changepage} %for indenting environments
\usepackage{csquotes} % example: \textcquote{}
\usepackage{blkarray}
\setcounter{MaxMatrixCols}{20} % default for pmatrix is 10!!
\usepackage{ytableau}
\usepackage{array} %needed to increase the vertical length in a tabular

% drawing
\usepackage{tikz,tikz-cd}
\usetikzlibrary{shapes.misc, patterns, matrix, calc, intersections,positioning,arrows.meta}
\usepackage{graphics,graphicx}
\usepackage{float} % provides enhanced control and customization options for floating objects such as figures and tables

% colors
\usepackage{xcolor}
\definecolor{darkcandyapplered}{rgb}{0.64, 0.0, 0.0}
\definecolor{midnightblue}{rgb}{0.1, 0.1, 0.44}
\definecolor{mylightblue}{HTML}{336699}
\definecolor{burntorange}{rgb}{0.8, 0.33, 0.0}
\definecolor{iceberg}{rgb}{0.44, 0.65, 0.82}
\definecolor{applegreen}{rgb}{0.55, 0.71, 0.0}
\definecolor{canaryyellow}{rgb}{1.0, 0.94, 0.0}
\definecolor{lightgray}{rgb}{0.83, 0.83, 0.83}
\definecolor{vpurple}{HTML}{440154}

% hrefs
\usepackage{hyperref}
\usepackage[noabbrev,capitalize]{cleveref}
\hypersetup{
    pdftoolbar=true,        
    pdfmenubar=true,        
    pdffitwindow=false,     
    pdfstartview={FitH},  % fits the width of the page to the window
    pdftitle={},
    pdfauthor={},
    pdfsubject={},
    pdfkeywords={},
    pdfnewwindow=true,  % links in new window
    colorlinks=true,  % false: boxed links; true: colored links
    linkcolor=darkcandyapplered,   % color of internal links
    citecolor=midnightblue,  % color of links to bibliography
    urlcolor=cyan,  % color of external links
    linktocpage=true  % changes the links from the section body to the page number
    }

% geometry
\textwidth=16cm 
\textheight=21cm 
\hoffset=-55pt 
\footskip=25pt

% thm envs (you might need to change the path)
\theoremstyle{definition}
\newtheorem{exercise}{Άσκηση}
\newtheorem{solution}{Λύση}
\newtheorem*{remark}{Παρατήρηση}
\newtheorem*{example}{Παράδειγμα}

% math ops 
% fields
\renewcommand{\AA}{\mathbbmss{A}} 
\newcommand{\NN}{\mathbbmss{N}} 
\newcommand{\ZZ}{\mathbbmss{Z}} 
\newcommand{\QQ}{\mathbbmss{Q}} 
\newcommand{\RR}{\mathbbmss{R}} 
\newcommand{\CC}{\mathbbmss{C}} 
\newcommand{\KK}{\mathbbmss{K}} 
\newcommand{\FF}{\mathbbmss{F}} 
\newcommand{\HH}{\mathbbmss{H}} 

% symmetric group
\newcommand{\fS}{\mathfrak{S}}  
\newcommand{\fp}{\mathfrak{p}}  
\newcommand{\fm}{\mathfrak{m}}  

% calligraphic 
\newcommand{\aA}{\mathcal{A}} 
\newcommand{\bB}{\mathcal{B}}
\newcommand{\cC}{\mathcal{C}}
\newcommand{\dD}{\mathcal{D}}
\newcommand{\eE}{\mathcal{E}}
\newcommand{\fF}{\mathcal{F}}
\newcommand{\hH}{\mathcal{H}}
\newcommand{\iI}{\mathcal{I}}
\newcommand{\lL}{\mathcal{L}}
\newcommand{\nN}{\mathcal{N}}
\newcommand{\oO}{\mathcal{O}}
\newcommand{\pP}{\mathcal{P}}
\newcommand{\sS}{\mathcal{S}}
\newcommand{\mM}{\mathcal{M}}
\newcommand{\uU}{\mathcal{U}}
\newcommand{\vV}{\mathcal{V}}

% bold
\newcommand{\bfa}{\mathbf{a}}
\newcommand{\bfe}{\mathbf{e}}
\newcommand{\bfF}{\pmb{F}}
\newcommand{\bfR}{\pmb{R}}
\newcommand{\bfv}{\mathbf{v}}
%\newcommand{\bfx}{\bm{x}}
%\newcommand{\bfx}{\mathbf{x}} 
\newcommand{\bfx}{\pmb{x}}
\newcommand{\bfX}{\pmb{X}}
\newcommand{\bfy}{\pmb{y}}
\newcommand{\bfz}{\pmb{z}}

% roman
\newcommand{\rmA}{\mathrm{A}}
\newcommand{\rmB}{\mathrm{B}}
\newcommand{\rmC}{\mathrm{C}}
\newcommand{\rmD}{\mathrm{D}} 
\newcommand{\rmH}{\mathrm{H}} 
\newcommand{\rmh}{\mathrm{h}} 
\newcommand{\rmI}{\mathrm{I}} 
\newcommand{\rmJ}{\mathrm{J}} 
\newcommand{\rmK}{\mathrm{K}}
\newcommand{\rmM}{\mathrm{M}}
\newcommand{\rmP}{\mathrm{P}}  
\newcommand{\rmp}{\mathrm{p}}  
\newcommand{\rmQ}{\mathrm{Q}}  
\newcommand{\rmR}{\mathrm{R}}
\newcommand{\rmS}{\mathrm{S}}
\newcommand{\rmT}{\mathrm{T}}
\newcommand{\rmU}{\mathrm{U}}
\newcommand{\rmV}{\mathrm{V}}
\newcommand{\rmY}{\mathrm{Y}}
\newcommand{\rmZ}{\mathrm{Z}}
\newcommand{\rmz}{\mathrm{z}}
\newcommand{\SL}{\mathrm{SL}}

% arrows and symbols 
\renewcommand{\to}{\rightarrow}
\newcommand{\toto}{\longrightarrow}
\newcommand{\mapstoto}{\longmapsto}
\newcommand{\then}{\Rightarrow}
\newcommand{\IFF}{\Leftrightarrow}
\newcommand{\tl}{\tilde}
\newcommand{\wtl}{\widetilde}
\newcommand{\ol}{\overline}
\newcommand{\ul}{\underline}
\newcommand{\oldemptyset}{\emptyset}
\renewcommand{\emptyset}{\varnothing}
\DeclareMathSymbol{\Arg}{\mathbin}{AMSa}{"39} % for arguments 
\newcommand{\onto}{\ensuremath{\twoheadrightarrow}}
\newcommand{\tle}{\trianglelefteq}
\newcommand{\tge}{\trianglerighteq}

% custom ops
\newcommand{\rmKer}{\mathrm{Ker}}
\newcommand{\rmIm}{\mathrm{Im}}
\newcommand{\lcm}{\mathrm{lcm}}
\newcommand{\GL}{\mathrm{GL}}
\newcommand{\ev}{\mathrm{ev}}
\newcommand{\End}{\mathrm{End}}
\newcommand{\rmid}{\mathrm{id}}
\newcommand{\rmspan}{\mathrm{span}}
\newcommand{\Spec}{\mathrm{Spec}}
\newcommand{\mSpec}{\mathrm{mSpec}}
\newcommand{\rad}{\mathrm{rad}}
\newcommand{\Rad}{\mathrm{Rad}}
\newcommand{\Jac}{\mathrm{Jac}}
\newcommand{\tor}{\mathrm{tor}}
\newcommand{\Ann}{\mathrm{Ann}}
\newcommand{\Hom}{\mathrm{Hom}}
\newcommand{\Hilb}{\mathrm{Hilb}}
\newcommand{\Ass}{\mathrm{Ass}}

% absolute value symbol
\usepackage{mathtools} 
\DeclarePairedDelimiter\abs{\lvert}{\rvert}%
\DeclarePairedDelimiter\norm{\lVert}{\rVert}%
\makeatletter
\let\oldabs\abs
\def\abs{\@ifstar{\oldabs}{\oldabs*}}

% tensor symbol
\newcommand{\tensor}[1]{%
  \mathbin{\mathop{\otimes}\limits_{#1}}%
}

% permutation cycle notation
\ExplSyntaxOn
\NewDocumentCommand{\cycle}{ O{\;} m }
 {
  (
  \alec_cycle:nn { #1 } { #2 }
  )
 }

\seq_new:N \l_alec_cycle_seq
\cs_new_protected:Npn \alec_cycle:nn #1 #2
 {
  \seq_set_split:Nnn \l_alec_cycle_seq { , } { #2 }
  \seq_use:Nn \l_alec_cycle_seq { #1 }
 }
\ExplSyntaxOff

% setminus symbol
\newcommand{\mysetminusD}{\hbox{\tikz{\draw[line width=0.6pt,line cap=round] (3pt,0) -- (0,6pt);}}}
\newcommand{\mysetminusT}{\mysetminusD}
\newcommand{\mysetminusS}{\hbox{\tikz{\draw[line width=0.45pt,line cap=round] (2pt,0) -- (0,4pt);}}}
\newcommand{\mysetminusSS}{\hbox{\tikz{\draw[line width=0.4pt,line cap=round] (1.5pt,0) -- (0,3pt);}}}
\newcommand{\sm}{\mathbin{\mathchoice{\mysetminusD}{\mysetminusT}{\mysetminusS}{\mysetminusSS}}}

% miscellaneous commands
\newcommand{\defn}[1]{{\color{mylightblue}{#1}}}
\newcommand{\toDo}{{\bf\color{red} TODO}}
\newcommand{\toCite}{{\bf\color{green} CITE}}
\newcommand*{\vertbar}{\rule[-1ex]{0.5pt}{2.5ex}} % for matrices with column vectors
\newcommand*{\horzbar}{\rule[.5ex]{2.5ex}{0.5pt}} % for matrices with row vectors
\newcommand{\myblue}[1]{{\color{iceberg}{#1}}}
\newcommand{\myorange}[1]{{\color{burntorange}{#1}}}
\newcommand{\mygreen}[1]{{\color{applegreen}{#1}}}
\newcommand{\myred}[1]{{\color{darkcandyapplered}{#1}}}

% ferrer's diagram
\newcommand{\fdiagram}[1]{
    \begin{tikzpicture}[scale=.7]
        \fill foreach \Z [count=\Y] in {#1}
        {foreach \X in {1,...,\Z} 
        {(\X,-\Y) circle[radius=3pt]}};
    \end{tikzpicture}
}

%
\usepackage{pgfplots}
\pgfplotsset{compat=1.18}

%
\makeatletter
\def\mathcolor#1#{\@mathcolor{#1}}
\def\@mathcolor#1#2#3{%
  \protect\leavevmode
  \begingroup
    \color#1{#2}#3%
  \endgroup
}
\makeatother

% 
\newenvironment{nouppercase}{%
  \let\uppercase\relax%
  \renewcommand{\uppercasenonmath}[1]{}}{}

% titlepage
\title{226: Θεωρία Δακτυλίων και Modules \\ Ασκήσεις με Ενδεικτικές Λύσεις}
\author[Β.~Δ. Μουστακας]{Βασίλης Διονύσης Μουστάκας \\ Πανεπιστήμιο Κρήτης}
\date{20 Μαΐου 2026}
% \urladdr{\href{https://sites.google.com/view/vasmous}{https://sites.google.com/view/vasmous}}

\begin{document}

\begingroup
\def\uppercasenonmath#1{} % this disables uppercase title
\let\MakeUppercase\relax % this disables uppercase authors
\maketitle
\endgroup

\setcounter{exercise}{48}
% \setcounter{theorem}{5}
\setcounter{equation}{10}
\renewcommand{\thesubsection}{\thesection.\arabic{subsection}}

Σε ότι ακολουθεί, με $F$ συμβολίζουμε ένα (άπειρο) σώμα.

\begin{proof}[Λύση του Ερωτήματος (2) της Άσκησης 48]
  Έστω $\fm = (x, y, z)$. Τότε, 
  \[
  F \cong F[x,y,z]/\fm
  \]
  και γι αυτό το $\iI(V)/\fm\iI(V)$ είναι διανυσματικός χώρος πάνω από το $F$, όπως στην απόδειξη του Θεωρήματος 5.6. Ισχυριζόμαστε ότι 
  \[
  \dim_F\left(
    \iI(V)/\fm\iI(V)
  \right) \ge 3.
  \]
  Πράγματι, τα στοιχεία 
  \[
  xy + \fm\iI(V), \ 
  xz + \fm\iI(V), \  
  yz + \fm\iI(V)
  \]
  είναι γραμμικά ανεξάρτητα, διότι διαφορετικά θα είχαμε 
  \[
  (axy + bxz + cyz) + \fm\iI(V) = 0 + \fm\iI(V)
  \]
  για κάποια $a, b, c \in F$, που θα σήμαινε ότι 
  \[
  axy + bxz + cyz \in \fm\iI(V),
  \]
  πράγμα αδύνατο, καθώς το 
  \[
  \fm\iI(V) = (x,y,z)(xy, xz, yz) = \{x^2y, x^2z, xy^2, xz^2, y^2z, yz^2, xyz, \dots\}
  \]
  δεν περιέχει όρους βαθμού $\le 2$.

  Όμως, αν $\{f_1, f_2, \dots, f_n\}$ είναι ένα σύνολο γεννητόρων του $\iI(V)$, τότε το σύνολο $\{f_1+\fm\iI(V), f_2+\fm\iI(V), \dots, f_n+\fm\iI(V)\}$ παράγει τον διανυσματικό χώρο $\iI(V)/\fm\iI(V)$ (γιατί;). Επειδή η διάσταση του τελευταίου είναι τουλάχιστον $3$, το ζητούμενο έπεται.
\end{proof}

% fulton 1.23
\begin{exercise}
  Να βρείτε μια διάσπαση της πολλαπλότητας 
  \[      
  V = \vV(\{y^4 - x^2, y^4 - x^2y^2 + xy^2 - x^3\})
  \]   
  σε ανάγωγες συνιστώσες.
\end{exercise}

\begin{proof}[Λύση]
  Αν $(a,b) \in V$, τότε 
  \begin{align*}
    b^4 - a^2 &= 0 \\
    b^4 - a^2b^2 + ab^2 - a^3 &= 0.
  \end{align*}
  Από την πρώτη σχέση έχουμε 
  \[
  0 = b^4 - a^2 = (b^2-a)(b^2+a)
  \]
  και γι αυτό είτε $b^2-a = 0$ ή $b^2 + a = 0$. Διακρίνουμε περιπτώσεις
  \begin{itemize}
    \item Αν $b^2 - a = 0$, τότε από την δεύτερη σχέση έχουμε 
    \begin{align*}
    0 &= b^4 - a^2b^2 + ab^2 - a^3 \\ 
    &= a^2 - a^2a + aa - a^3 \\
    &= a^2 - a^3 + a^2 - a^3\\
    &= 2a^2(1-a)
    \end{align*}
    και γι αυτό είτε $a^2 = 0$ που σημαίνει ότι $a=0$ ή $a = 1$. Συνεπώς, στην περίπτωση αυτή έχουμε τις επιλογές 
    \[
    (a,b) \in \{
      (0,0), (1,1), (1,-1)
    \}.
    \]
    \item Αν $b^2 + a = 0$, τότε η δεύτερη σχέση ικανοποιείται ταυτοτικά (γιατί;). Συνεπώς, στην περίπτωση αυτή ολόκληρη η πολλαπλότητα $\vV(y^2+x)$ περιέχεται στην $V$.
  \end{itemize}
  Επειδή $(0,0) \in \vV(y^2+x)$, έχουμε 
  \[
  V = \vV(y^2+x) \cup \vV(\{x-1,y-1\}) \cup \vV(\{x-1,y+1\}),
  \]
  από το Παράδειγμα 9.3 (δ). Αυτή είναι η διάσπαση σε ανάγωγες συνιστώσες (γιατί;).
\end{proof}

% DF example on page 660
\begin{exercise}
  Αν $V = \vV(x^3-y^2) \subseteq \AA^2$, τότε να δείξετε ότι $\iI(V) = (x^3-y^2)$.
\end{exercise}

\begin{proof}[Λύση]
  Αρχικά παρατηρούμε ότι 
  \[
  V = \{(t^2, t^3) : t \in F\}.
  \]
  Πράγματι, αν $(a, b) \in V$, τότε $a^3 = b^2$. Αν $a \neq 0$, τότε και $b \neq 0$ και θέτοντας $t = b/a \in F$ έχουμε
  \begin{align*}
    t^2 &= \frac{b^2}{a^2} = \frac{a^3}{a^2} = a \\
    t^3 &= \frac{b^3}{a^3} = \frac{b^3}{b^2} = b.
  \end{align*}
  Για τον αντίστροφο εγκλεισμό, αρκεί να παρατηρήσουμε ότι 
  \[
  \ev_{t^2, t^3}(x^3-y^2) = t^6 - t^6 = 0.
  \]

  Έπειτα, παρατηρούμε ότι $\iI(V) = (x^3-y^2)$. Πράγματι, για κάθε $f(x,y) \in F[x,y]$ από την γενικευμένη εκδοχή του Θεωρήματος 3.3 για $R = F[x]$ και $g(\mathcolor{gray}{x},y) = x^3-y^2$ έπεται ότι 
  \[
  f(\mathcolor{gray}{x},y) = q(\mathcolor{gray}{x},y)(\mathcolor{gray}{x^3}-y^2) + r(\mathcolor{gray}{x},y),
  \]
  όπου $q, r \in R[y]$ με $\deg(r) <2$. Συνεπώς,
  \[
  r(\mathcolor{gray}{x},y) = \mathcolor{gray}{f_0(x)} + \mathcolor{gray}{f_1(x)}y
  \]
  για κάποια $f_0(x), f_1(x) \in F[x]$. Αν $f(x,y) \in \iI(V)$, τότε $f(t^2,t^3) = 0$, για κάθε $t \in F$ και γι αυτό 
  \[
  f_0(t^2) + f_1(t^2)t^3 = 0,
  \]
  για κάθε $t \in F$. Αυτό σημαίνει ότι το πολυώνυμο μιας μεταβλητής 
  \begin{align*}
  r(X^2, X^3) 
  &= f_0(X^2) + f_1(X^2)X^3 \\
  &= \left(
    a_0 + a_1X^2 + \cdots + a_kX^{2k} 
  \right) + \left(
    b_0X^3 + b_1X^5 + \cdots + b_{\ell}X^{2\ell+3} 
  \right) \\
  &=0
  \end{align*}
  Οι συντελεστές των όρων άρτιου (αντ. περιττού) βαθμού είναι οι συντελεστές των $f_0(x)$ και $f_1(x)$, αντίστοιχα και γι αυτό $f_0(x) = f_1(x) = 0$. Άρα, 
  \[
  f(x,y) = q(x,y)(x^3-y^2),
  \]
  και γι αυτό $f \in (x^3-y^2)$, όπως θέλαμε.
\end{proof}

%Cox Prop 4.2.9 + DF 15.1.19
\begin{exercise}{\cite[Άσκηση~15.1.19]{DF04}}
  Έστω $f \in F[x_1, x_2, \dots, x_n]$ και $I = (f)$. Αν 
  \[
  f = c f_1^{k_1}f_2^{k_2}\cdots f_m^{k_m}
  \]
  είναι η παραγοντοποίηση του $f$ σε ανάγωγα στοιχεία του $F[x_1, \dots, x_n]$, τότε να δείξετε ότι 
  \[
  \sqrt{I} = \left(
    f_1f_2\cdots{f_m}
  \right).
  \]
  Συμπεράνετε τα εξής:
  \begin{itemize}
    \item Το $I$ είναι ριζικό ιδεώδες αν και μόνο αν το $f$ είναι ελεύθερο τετραγώνων
    \item (\cite[Άσκηση~15.2.15]{DF04}) Αν το $F$ είναι αλγεβρικά κλειστό, τότε η υπερεπιφάνεια $\vV(f)$ είναι ανάγωγη αν και μόνο αν το $f$ είναι ανάγωγο.
  \end{itemize}
\end{exercise}

\begin{proof}[Λύση]
  Οι δύο τελευταίοι ισχυρισμοί αφήνονται ως ασκήσεις. Για τον πρώτο ισχυρισμό, εργαζόμαστε ως εξής:
  Για την κατεύθυνση \textquote{$\supseteq$}, αν $k = \max\{k_1, k_2, \dots, k_m\}$, τότε 
  \[
  \left(
    f_1f_2\cdots f_m
  \right)^k = \underbrace{f_1^{k-k_1}f_2^{k-k_2}\cdots f_m^{k-k_m}}_{\in \ F[x_1, x_2, \dots, x_n]} f \ \in \ I.
  \]
  Οπότε, $f_1f_2\cdots f_m \in \sqrt{I}$ και γι αυτό 
  \[
  \left(
    f_1f_2\cdots{f_m}
  \right) \subseteq \sqrt{I}.
  \]

  Για την αντίστροφη κατεύθυνση, έστω $g \in \sqrt{I}$. Τότε $g^k \in I$, για κάποιο $k \ge 1$, και γι αυτό $g^k = hf$, για κάποιο 
  \[
  g^k = h f_1^{k_1}f_2^{k_2}\cdots f_m^{k_m}
  \]
  για κάποιο $h \in F[x_1, x_2, \dots, x_n]$. Από το Θεώρημα 3.6 γνωρίζουμε ότι ο πολυωνυμικός δακτύλιος είναι περιοχή μονοσήμαντης παραγοντοποίησης και γι αυτό $f_i \mid g^k$, για κάθε $1 \le i \le m$. Από το Λήμμα 2.14 γνωρίζουμε ότι σε μια περιοχή μονοσήμαντης παραγοντοποίησης κάθε ανάγωγο στοιχείο είναι και πρώτο και γι αυτό $f_i \mid g$, για κάθε $1 \le i \le m$. Με άλλα λόγια, $g \in (f_i)$, για κάθε $1 \le i \le m$. Οπότε, 
  \[
  g \in \left(
    f_1f_2\cdots{f_m}
  \right)
  \]
  και το ζητούμενο έπεται.
\end{proof}

\begin{remark}
  Αν το $F$ δεν είναι αλγεβρικά κλειστό, τότε δεν ισχύει απαραίτητα. Για παράδειγμα, το πολυώνυμο $f(x,y) = y^2 + x^2(x-1)^2$ είναι ανάγωγο στο $\RR[x,y]$ (γιατί;), αλλά η αντίστοιχη υπερεπιφάνεια του $\AA_\RR^2$ δεν είναι.
\end{remark}

\begin{thebibliography}{99}
    \bibitem{DF04}
    D. S. Dummit και R. M. Foote, 
    \textit{Abstract algebra},
    third edition, Wiley, 2004. 
\end{thebibliography}
\end{document}